\section{Design}
\label{sec:design}

Chiyoda models evacuation as a feedback loop between physical state,
information state, and action. At each step, hazards evolve; agents observe
nearby exits and hazards; beacons, responders, gossip, and intervention
policies update beliefs; agents revise intentions; belief-weighted navigation
selects paths; social-force dynamics move agents; and telemetry records the
macroscopic outcome.

\subsection{Belief State}
\label{subsec:beliefs}

Each agent carries a belief vector over exits, hazards, and general danger.
Exit beliefs store existence probability, congestion estimate, freshness,
source credibility, and gossip hop count. Hazard beliefs store position,
severity estimate, radius estimate, freshness, credibility, and hop count.
Agents do not route through the ground truth environment directly; they route
through the subset of exits and hazards represented in their beliefs.

\subsection{Communication Channels}
\label{subsec:channels}

The baseline ITED channels are direct observation, beacon broadcast, and
agent-to-agent gossip. Observation is range-limited by visibility and
physiological impairment. Beacons provide high-credibility exit knowledge
within a fixed radius. Gossip is local, probabilistic, and distorted by hop
count and panic state.

\subsection{Intervention Policies}
\label{subsec:policies}

The new contribution is an intervention policy interface. Every policy decides
whether to fire, selects one or more spatial targets, builds a message, applies
it to recipients inside a radius, and emits telemetry. Chiyoda currently
implements seven policies:

\begin{itemize}
\item \textbf{Static beacon}: periodic local messages from existing signage
or PA beacon locations.
\item \textbf{Global broadcast}: station-wide high-reach communication.
\item \textbf{Responder relay}: mobile high-credibility messages emitted by
released responders.
\item \textbf{Entropy-targeted}: targets agents with highest belief entropy.
\item \textbf{Density-aware}: targets dense local clusters.
\item \textbf{Exposure-aware}: targets agents under high current or cumulative
hazard pressure.
\item \textbf{Bottleneck-avoidance}: targets active bottleneck zones.
\end{itemize}

Figure~\ref{fig:policy-taxonomy} summarizes the policy space. Figure~\ref{fig:intervention-loop}
shows where the intervention hook sits in the simulation loop.

\begin{figure}[!t]
  \centering
  \includegraphics[width=\linewidth]{figures/policy-taxonomy.pdf}
  \caption{Information intervention policy taxonomy. Policies differ by
  source, target signal, spatial reach, and safety objective.}
  \label{fig:policy-taxonomy}
\end{figure}

\begin{figure}[!t]
  \centering
  \includegraphics[width=\linewidth]{figures/intervention-loop.pdf}
  \caption{Intervention hook in the ITED simulation loop. Policies read the
  current physical-information state, update recipient beliefs, and emit
  before/after telemetry before navigation.}
  \label{fig:intervention-loop}
\end{figure}

\subsection{Information-Safety Metrics}
\label{subsec:metrics}

The core metric is information-safety efficiency:

\[
\eta_{safe} =
\frac{\Delta H^{+} + \Delta A^{+}}
     {1 + P_{exposure} + P_{queue}},
\]

where $\Delta H^{+}$ is positive entropy reduction, $\Delta A^{+}$ is positive
belief-accuracy gain, $P_{exposure}$ is recipient-weighted hazard-load
pressure, and $P_{queue}$ is bottleneck pressure observed at intervention
time. Chiyoda also records intervention reach, exit imbalance, and a harmful
convergence index that rises when entropy reduction coincides with concentrated
exit usage, queueing, and exposure.
